% GENERATED FILE. Edit canonical lessons, then run npm run generate:books.
% canonical-source-hash: fnv1a64:ed4bb4ac66388f2f
% canonical-lessons: JA-C447-memai, JA-C447-furafura, JA-C447-darui, JA-C447-guttari, JA-C447-nebusoku

\chapter{Worn Out}
\label{ch:ja-worn-out}
\hlchaptermodality{447}

\begin{chapteropening}
\textbf{By the end of this chapter:} \emph{I can say I feel dizzy, sluggish or worn out.}

Complete the last lesson of chapter 447: I can say I feel dizzy, sluggish or worn out.
\end{chapteropening}

\section[\texorpdfstring{\ja{めまい} (memai)}{めまい (memai)}]{\ja{めまい} (memai) — dizziness}

\label{lesson:JA-C447-memai}



\begin{quote}
Before the new one: say the Japanese for gargling, then the Japanese for physical condition.
\end{quote}

\subsection*{You'll want to know: \ja{めまい}}
\textbf{\ja{めまい}} — \emph{memai} — \textquotedblleft{}dizziness\textquotedblright{}.

\begin{quote}
\textbf{\ja{あさ} \ja{おきた} \ja{とき、めまいが} \ja{しました。}} — \emph{asa okita toki, memai ga shimashita} — \textquotedblleft{}When I got up this morning, I felt dizzy.\textquotedblright{}
\end{quote}

\begin{cousinweb}
\textbf{\ja{めまい}} (\emph{memai}) is a native Japanese word (wago), inherited rather than borrowed from Chinese. JMdict marks it as usually written in kana and lists five kanji spellings for it, so no single spelling shows where it comes from (JMdict).
\end{cousinweb}

\begin{grammarlens}[title={What you've built}]
One more word for describing what you lived and felt.
\end{grammarlens}

\subsection*{Your turn}
\begin{itemize}
\raggedright
  \item \emph{Say it:} \emph{memai}
  \item \emph{Say it:} \emph{memai}, once more
  \item \emph{Say it:} \emph{asa okita toki, memai ga shimashita} — When I got up this morning, I felt dizzy.
  \item \emph{Recall:} say the Japanese for gargling, then the Japanese for physical condition, then say \emph{memai} again
\end{itemize}

\begin{tcolorbox}[breakable,colback=teal!4,colframe=teal!35!black,title={Before you move on}]
What does \ja{めまい} mean? (\textquotedblleft{}Dizziness\textquotedblright{}.) Say it once more.
\end{tcolorbox}
\section[\texorpdfstring{\ja{ふらふら} (furafura)}{ふらふら (furafura)}]{\ja{ふらふら} (furafura) — unsteadily; dizzily}

\label{lesson:JA-C447-furafura}



\begin{quote}
Before the new one: say the Japanese for physical condition, then the Japanese for dizziness.
\end{quote}

\subsection*{You'll want to know: \ja{ふらふら}}
\textbf{\ja{ふらふら}} — \emph{furafura} — \textquotedblleft{}unsteadily; dizzily\textquotedblright{}.

\begin{quote}
\textbf{\ja{ねつで} \ja{ふらふらします。}} — \emph{netsu de furafura shimasu} — \textquotedblleft{}The fever makes me unsteady on my feet.\textquotedblright{}
\end{quote}

\begin{cousinweb}
\textbf{\ja{ふらふら}} (\emph{furafura}) is a native Japanese mimetic word, inherited rather than borrowed from Chinese, and its older origin is not settled (JMdict).
\end{cousinweb}

\begin{grammarlens}[title={What you've built}]
One more word for describing what you lived and felt.
\end{grammarlens}

\subsection*{Your turn}
\begin{itemize}
\raggedright
  \item \emph{Say it:} \emph{furafura}
  \item \emph{Say it:} \emph{furafura}, once more
  \item \emph{Say it:} \emph{netsu de furafura shimasu} — The fever makes me unsteady on my feet.
  \item \emph{Recall:} say the Japanese for physical condition, then the Japanese for dizziness, then say \emph{furafura} again
\end{itemize}

\begin{tcolorbox}[breakable,colback=teal!4,colframe=teal!35!black,title={Before you move on}]
What does \ja{ふらふら} mean? (\textquotedblleft{}Unsteadily; dizzily\textquotedblright{}.) Say it once more.
\end{tcolorbox}
\section[\texorpdfstring{\ja{だるい} (darui)}{だるい (darui)}]{\ja{だるい} (darui) — sluggish; listless}

\label{lesson:JA-C447-darui}



\begin{quote}
Before the new one: say the Japanese for dizziness, then the Japanese for unsteadily.
\end{quote}

\subsection*{You'll want to know: \ja{だるい}}
\textbf{\ja{だるい}} — \emph{darui} — \textquotedblleft{}sluggish; listless\textquotedblright{}.

\begin{quote}
\textbf{\ja{かぜで} \ja{からだが} \ja{だるいです。}} — \emph{kaze de karada ga darui desu} — \textquotedblleft{}My cold makes my body feel sluggish.\textquotedblright{}
\end{quote}

An i-adjective.

\begin{cousinweb}
\textbf{\ja{だるい}} (\emph{darui}) is a native Japanese word (wago), inherited rather than borrowed from Chinese; its kanji spelling was attached to it later (JMdict).
\end{cousinweb}

\begin{grammarlens}[title={What you've built}]
One more word for describing what you lived and felt.
\end{grammarlens}

\subsection*{Your turn}
\begin{itemize}
\raggedright
  \item \emph{Say it:} \emph{darui}
  \item \emph{Say it:} \emph{darui}, once more
  \item \emph{Say it:} \emph{kaze de karada ga darui desu} — My cold makes my body feel sluggish.
  \item \emph{Recall:} say the Japanese for dizziness, then the Japanese for unsteadily, then say \emph{darui} again
\end{itemize}

\begin{tcolorbox}[breakable,colback=teal!4,colframe=teal!35!black,title={Before you move on}]
What does \ja{だるい} mean? (\textquotedblleft{}Sluggish; listless\textquotedblright{}.) Say it once more.
\end{tcolorbox}
\section[\texorpdfstring{\ja{ぐったり} (guttari)}{ぐったり (guttari)}]{\ja{ぐったり} (guttari) — worn out; limp}

\label{lesson:JA-C447-guttari}



\begin{quote}
Before the new one: say the Japanese for unsteadily, then the Japanese for sluggish.
\end{quote}

\subsection*{You'll want to know: \ja{ぐったり}}
\textbf{\ja{ぐったり}} — \emph{guttari} — \textquotedblleft{}worn out; limp\textquotedblright{}.

\begin{quote}
\textbf{\ja{なつに} \ja{ながく} \ja{あるいて、ぐったりしました。}} — \emph{natsu ni nagaku aruite, guttari shimashita} — \textquotedblleft{}I walked a long way in summer and was worn out.\textquotedblright{}
\end{quote}

\begin{cousinweb}
\textbf{\ja{ぐったり}} (\emph{guttari}) is a native Japanese mimetic word, inherited rather than borrowed from Chinese, and its older origin is not settled (JMdict).
\end{cousinweb}

\begin{grammarlens}[title={What you've built}]
One more word for describing what you lived and felt.
\end{grammarlens}

\subsection*{Your turn}
\begin{itemize}
\raggedright
  \item \emph{Say it:} \emph{guttari}
  \item \emph{Say it:} \emph{guttari}, once more
  \item \emph{Say it:} \emph{natsu ni nagaku aruite, guttari shimashita} — I walked a long way in summer and was worn out.
  \item \emph{Recall:} say the Japanese for unsteadily, then the Japanese for sluggish, then say \emph{guttari} again
\end{itemize}

\begin{tcolorbox}[breakable,colback=teal!4,colframe=teal!35!black,title={Before you move on}]
What does \ja{ぐったり} mean? (\textquotedblleft{}Worn out; limp\textquotedblright{}.) Say it once more.
\end{tcolorbox}
\section[\texorpdfstring{\ja{ねぶそく} (nebusoku)}{ねぶそく (nebusoku)}]{\ja{ねぶそく} (nebusoku) — lack of sleep}

\label{lesson:JA-C447-nebusoku}



\begin{quote}
Before the new one: say the Japanese for sluggish, then the Japanese for worn out.
\end{quote}

\subsection*{You'll want to know: \ja{ねぶそく}}
\textbf{\ja{ねぶそく}} — \emph{nebusoku} — \textquotedblleft{}lack of sleep\textquotedblright{}.

\begin{quote}
\textbf{\ja{ねぶそくで} \ja{あたまが} \ja{いたいです。}} — \emph{nebusoku de atama ga itai desu} — \textquotedblleft{}My head aches from lack of sleep.\textquotedblright{}
\end{quote}

\begin{cousinweb}
\textbf{\ja{ねぶそく}} (\emph{nebusoku}) is a mixed word: native \emph{ne}, \textquotedblleft{}sleeping\textquotedblright{}, from \emph{neru}, joined to Sino-Japanese \emph{busoku}, \textquotedblleft{}shortage\textquotedblright{}, from \emph{bu}, \textquotedblleft{}negative\textquotedblright{} (Middle Chinese \emph{pjuw}), and \emph{soku}, \textquotedblleft{}be sufficient\textquotedblright{} (\emph{tsjowk}), readings borrowed from Chinese (JMdict, KANJIDIC2; Middle Chinese in Baxter's transcription).
\end{cousinweb}

\begin{grammarlens}[title={What you've built}]
One more word for describing what you lived and felt.
\end{grammarlens}

\subsection*{Your turn}
\begin{itemize}
\raggedright
  \item \emph{Say it:} \emph{nebusoku}
  \item \emph{Say it:} \emph{nebusoku}, once more
  \item \emph{Say it:} \emph{nebusoku de atama ga itai desu} — My head aches from lack of sleep.
  \item \emph{Recall:} say the Japanese for sluggish, then the Japanese for worn out, then say \emph{nebusoku} again
\end{itemize}

\begin{tcolorbox}[breakable,colback=teal!4,colframe=teal!35!black,title={Before you move on}]
What does \ja{ねぶそく} mean? (\textquotedblleft{}Lack of sleep\textquotedblright{}.) Say it once more.
\end{tcolorbox}
